% !TEX TS-program = lualatex

На \refimage{high-level-design} изображено представление программного решения на высоком уровне проектирования архитектуры.

\begin{enumerate}
    \item Модуль ядра (game\_core) является подключаемой библиотекой и выполняет роль основного контроллера для управления ходом игры и её состоянием. Он обеспечивает безопасный доступ к моделям представления, содержит описания правил для больших и малых вариантов игры <<Нарды>>, а также предоставляет базовую визуализацию игрового поля.
    \item Модуль ядра машинного обучения (model\_core) представляет собой подключаемую библиотеку, отвечающую за реализацию и интеграцию алгоритмов машинного обучения в игровой процесс. Этот модуль содержит описания моделей, а также инструменты, позволяющие загружать и выгружать их из оперативной памяти. Он взаимодействует с модулем ядра, получая доступ к текущему состоянию игрового поля, истории ходов и другим данным, необходимым для обучения и принятия решений.
    \item Модуль для тренировки интеллектуальных агентов (model\_trainer) представляет из себя инструмент командной строки, предназначенный для обучения и валидации моделей, используемых в игре <<Нарды>>. Он предоставляет функциональность для автоматического обучения моделей через спарринг, настройки гиперпараметров, а также мониторинга метрик обучения.
    \item Модуль сервиса, обеспечивающего доступность моделей машинного обучения (model\_set\_server), представляет серверный компонент, предназначенный для хранения, управления и предоставления доступа к обученным моделям. Он реализует интерфейс для загрузки, выгрузки и версионирования моделей, а также обеспечивает авторизованный доступ к ним.
    \item Модуль серверной компоненты игры (game\_server) представляет собой центральный элемент сетевой инфраструктуры, отвечающий за управление игровыми сессиями и обработку ходов игроков. Он обеспечивает взаимодействие между пользовательскими интерфейсами, модулем ядра и сервисом предоставления доступа к моделям машинного обучения.
    \item Модуль клиентской компоненты игры (game\_client), с которым взаимодействует игрок. Он управляет пользовательским интерфейсом для игры <<Нарды>>, включая визуализацию игрового поля, бросков костей, ходов и отображение состояний игры. Он отправляет ходы и действия игрока на игровой сервер и получает обновления о состоянии игры и действиях искусственного интеллекта, также управляя проверкой ввода данных игроком и визуальной обратной связью.
    \item Модуль для выгрузки натренированных интеллектуальных агентов (model\_uploader) представляет собой небольшой вспомогательный инструмент командной строки, предназначенный для передачи обученных моделей в сервис хранения и управления моделями. Он упрощает процесс загрузки, автоматически решая вопросы авторизации и разрешения конфликтов от имени пользователя.
\end{enumerate}

\image[width=0.8\textwidth]{Высокоуровневый дизайн архитектуры программного продукта}{high-level-design}{high-level-design}
